%! Unit Opener: Functions as Models — Unit 2, Lesson 2.0 — Guided Notes (Answer Key)
\documentclass[10pt]{article}
\usepackage{atda-article}
\usepackage{atda-key}   % pulls in -boxes; do NOT also load -boxes
\newcommand{\TallMath}[1]{$\displaystyle #1\rule[-1.4em]{0pt}{3.2em}$}
\setlength{\workrowsep}{6pt}

\begin{document}

\pageheader{Unit 2, Lesson 2.0}{Guided Notes --- Answer Key}
% NAMESTRIP: no \namedateperiod here — mirrors the blank. Teacher-only prose goes
% in the lesson plan as \begin{teachernote}[Guided Notes], never in a key.

\vspace{0.15in}

\begin{objectivebox}
\small
By the end of this lesson, I will be able to\dots
\begin{enumerate}[leftmargin=1.5em, itemsep=2pt, topsep=3pt]
  \item read a graph in context and name what it tells me: the \textbf{domain} and \textbf{range},
        where the function is \textbf{increasing}, \textbf{decreasing}, or \textbf{constant}, and
        its highest and lowest points;
  \item recognize the three function families this unit studies --- \textbf{linear},
        \textbf{quadratic}, and \textbf{exponential} --- from a graph or a table, and name the
        \textbf{parent function} of each;
  \item use the algebra Unit 2 is built on: function notation, values a formula excludes, rate of
        change, and the intercepts of a line;
  \item name what each lesson of Unit 2 teaches and where that skill gets used later in this
        course.
\end{enumerate}
\end{objectivebox}

\vspace{0.10in}

\boxguard
\begin{vocabbox}
\small
Fill in each term as we build it together. These five words carry the whole unit.
\vocabans{Function}{A rule that pairs each input with exactly one output. We write $f(x)$ for the
output the rule gives the input $x$; the point $(x, f(x))$ is on its graph.}
\vocabans{Domain and range}{The domain is every input the function is allowed to take; the range
is every output it produces. A situation can shrink both.}
\vocabans{Increasing, decreasing, constant}{Read the graph left to right: increasing means the
outputs rise, decreasing means they fall, constant means they hold steady.}
\vocabans{Absolute maximum and minimum}{The highest and lowest points on the graph over the
\emph{whole} domain. Report both the value and where it happens.}
\vocabans{Parent function}{The simplest member of a family, before any shifting, stretching, or
flipping: $f(x)=x$, $g(x)=x^2$, $h(x)=2^x$.}
\end{vocabbox}

\vspace{0.10in}

\boxguard[12]
\begin{hookbox}
\small
\textbf{The Saturday hike.} Your friend hikes a mountain trail and their watch records one thing:
elevation, every minute, for six hours. At the end it hands you \emph{a picture} --- a graph. No
formula. No table. Just the curve.

You want to know how long they hiked, how high they got, when they stopped for lunch, and how
steep the way down was. \textbf{Can a picture answer all four?}

\ansline{Yes --- all four. Length is how far the curve runs, height is its top,}\\
\ansline{lunch is the flat stretch, and steepness is how sharply it falls.}\\
\end{hookbox}

\vspace{0.10in}

\boxguard[30]
\begin{notesbox}{1. One Graph, Every Question}
\small
Here is the watch's graph. $E$ is elevation in feet above sea level, $t$ hours after the hikers
left the trailhead.

\begin{center}
\begin{tikzpicture}
\begin{axis}[
    width=12.2cm, height=5.6cm,
    xlabel={$t$ (hours after leaving the trailhead)},
    ylabel={$E$ (feet above sea level)},
    xmin=0, xmax=6.4, ymin=0, ymax=2300,
    xtick={0,1,2,3,4,5,6}, ytick={0,400,800,1200,1600,2000},
    grid=both, grid style={linegray!45},
    axis lines=left,
    tick label style={font=\scriptsize}, label style={font=\scriptsize}]
  \addplot[royal, very thick, mark=*, mark size=1.6pt]
    coordinates {(0,800) (2,1600) (3,1600) (4,2000) (6,1000)};
\end{axis}
\end{tikzpicture}
\end{center}

\vspace{-2pt}
Every question below is answered by \emph{looking}. Fill in the name of the characteristic and
what you read off the graph.

\renewcommand{\arraystretch}{1.5}
\begin{tabularx}{\linewidth}{@{}p{4.9cm} p{3.3cm} X@{}}
\rowcolor{mist}
\textbf{The question} & \textbf{What it is called} & \textbf{What you read off} \\
\hline
How long were they on the trail? & \ans{domain} & \ans{$0 \le t \le 6$ hours} \\
What elevation did they start at? & \ans{$E$-intercept} & \ans{$800$ feet} \\
When were they resting? & \ans{constant interval} & \ans{from $t=2$ to $t=3$} \\
How high did they get, and when? & \ans{absolute maximum} & \ans{$2000$ ft at $t=4$} \\
What was their lowest elevation? & \ans{absolute minimum} & \ans{$800$ ft at $t=0$} \\
\end{tabularx}

\vspace{6pt}
The curve never touches the $t$-axis, so this function has no \ans{zeros} --- the hikers never
come down to sea level. A graph may have none, one, or many.

\vspace{4pt}
They moved downhill faster than they climbed. How can you tell that from the picture alone?

\ansline{The last piece of the curve is the steepest one on the graph:}\\
\ansline{it drops 1000 feet in 2 hours, while the climb rose 400 feet per hour.}\\
\end{notesbox}

\vspace{0.10in}

\boxguard[22]
\begin{notesbox}{2. Three Families, Three Shapes}
\small
Unit 2 studies three families of functions. The simplest member of a family is its
\textbf{parent function} --- everything else in the family is that parent moved, stretched, or
flipped.

\begin{center}
\renewcommand{\arraystretch}{1.1}
\begin{tabular}{@{}c@{\hspace{0.5cm}}c@{\hspace{0.5cm}}c@{}}
\begin{tikzpicture}
\begin{axis}[width=4.7cm, height=4.3cm, xmin=-3.2, xmax=3.2, ymin=-3.2, ymax=3.2,
    axis lines=middle, xtick={-2,2}, ytick={-2,2},
    tick label style={font=\tiny}, axis line style={-}]
  \addplot[royal, very thick, domain=-3:3, samples=2]{x};
\end{axis}
\end{tikzpicture}
&
\begin{tikzpicture}
\begin{axis}[width=4.7cm, height=4.3cm, xmin=-3.2, xmax=3.2, ymin=-1.2, ymax=5.2,
    axis lines=middle, xtick={-2,2}, ytick={2,4},
    tick label style={font=\tiny}, axis line style={-}]
  \addplot[royal, very thick, domain=-2.2:2.2, samples=60]{x^2};
\end{axis}
\end{tikzpicture}
&
\begin{tikzpicture}
\begin{axis}[width=4.7cm, height=4.3cm, xmin=-3.2, xmax=3.2, ymin=-1.2, ymax=5.2,
    axis lines=middle, xtick={-2,2}, ytick={2,4},
    tick label style={font=\tiny}, axis line style={-}]
  \addplot[royal, very thick, domain=-3:2.2, samples=60]{2^x};
\end{axis}
\end{tikzpicture}
\\[-2pt]
{\footnotesize $f(x)=x$} & {\footnotesize $g(x)=x^2$} & {\footnotesize $h(x)=2^x$} \\
\end{tabular}
\end{center}

\vspace{-4pt}
\renewcommand{\arraystretch}{1.5}
\begin{tabularx}{\linewidth}{@{}p{2.5cm} p{2.4cm} X@{}}
\rowcolor{mist}
\textbf{Family} & \textbf{Parent} & \textbf{What gives it away on a graph} \\
\hline
Linear & $f(x)=x$ & A straight line: the same \ans{change} in $y$ for every step in $x$. \\
Quadratic & $g(x)=x^2$ & A U-shaped curve with exactly one \ans{turning} point and a line of
    symmetry. \\
Exponential & $h(x)=2^x$ & A curve that \ans{multiplies} by the same factor each step and flattens
    toward a line it never touches. \\
\end{tabularx}

\vspace{6pt}
Every other member of a family is the parent function after a \ans{transformation} --- a shift, a
stretch, or a flip. Lesson 2.2 gives each one a name.
\end{notesbox}

\vspace{0.10in}

\boxguard[26]
\begin{notesbox}{3. The Unit 2 Roadmap}
\small
Mark the last column for yourself using your warm-up results: \textbf{C} = confident,
\textbf{R} = rusty, \textbf{N} = new to me. You will come back to this page before the unit test.

\renewcommand{\arraystretch}{1.45}
\begin{tabularx}{\linewidth}{@{}c p{2.9cm} X c@{}}
\rowcolor{mist}
\textbf{Lesson} & \textbf{Standard} & \textbf{What I will be able to do} & \textbf{C/R/N} \\
\hline
2.1 & \texttt{AFDA.AF.2a, e} & Use function notation, and find domain and range --- including the
      limits a context puts on them & \blank{0.9cm} \\
2.2 & \texttt{AFDA.AF.1a--b} & Name the parent function of a family and describe the
      transformation that produced a given graph & \blank{0.9cm} \\
2.3 & \texttt{AFDA.AF.1d, f} & Go from an equation to a graph and from a graph to an equation,
      verified on the TI-84 & \blank{0.9cm} \\
2.4 & \texttt{AFDA.AF.2b--d} & Read intercepts, zeros, maximums and minimums, and
      increasing/decreasing intervals off a graph & \blank{0.9cm} \\
2.5 & \texttt{AFDA.AF.2f--g} & Describe end behavior and identify horizontal and vertical
      asymptotes & \blank{0.9cm} \\
2.6 & \texttt{AFDA.AF.2} & Read and analyze piecewise-defined functions & \blank{0.9cm} \\
2.7 & \texttt{AFDA.AF.1c, e, g}, \texttt{AF.2h} & Decide which family models a situation, and
      compare models using graphs, tables, and equations & \blank{0.9cm} \\
\end{tabularx}

\vspace{4pt}
\textbf{My goal for this unit:} \ansline{student response --- accept any specific, checkable goal}
\end{notesbox}

\vspace{0.10in}

\boxguard
\begin{notesbox}{4. What Makes This Unit Different}
\small
\begin{itemize}[leftmargin=1.3em, itemsep=3pt, topsep=3pt]
  \item In Unit 1 you were handed an \ans{expression} and asked to rewrite it. In Unit 2 you are
        handed a \ans{graph} and asked what it \emph{says}. Almost every question this unit
        starts from a picture.
  \item The five questions you asked about the trail --- how long, how high, how low, when
        rising, where it starts --- are the same five you will ask of a normal curve in Unit 6 and
        of a sine wave in Unit 8. Learn the language once; spend it all year.
  \item Unit 1 is still underneath. An excluded value like the one in your warm-up becomes a
        \textbf{restriction on the domain} in Lesson 2.1, and in Lesson 2.5 it becomes a vertical
        \textbf{asymptote} --- a line the graph runs alongside forever without touching.
\end{itemize}
\end{notesbox}

\vspace{0.10in}

\boxguard[30]
\begin{practicebox}
\small
The student council sells wristbands at the fall festival. If it charges $p$ dollars per
wristband, its revenue in dollars is graphed below for prices from \$0 to \$12.

\begin{center}
\begin{tikzpicture}
\begin{axis}[
    width=10.6cm, height=5.2cm,
    xlabel={$p$ (price per wristband, dollars)}, ylabel={$R$ (revenue, dollars)},
    xmin=0, xmax=12.8, ymin=0, ymax=210,
    xtick={0,2,4,6,8,10,12}, ytick={0,60,120,180},
    grid=both, grid style={linegray!45}, axis lines=left,
    tick label style={font=\scriptsize}, label style={font=\scriptsize}]
  \addplot[royal, very thick, domain=0:12, samples=80]{-5*x^2+60*x};
\end{axis}
\end{tikzpicture}
\end{center}

\begin{enumerate}[leftmargin=1.6em, itemsep=6pt, topsep=2pt]
  \item Where does the graph cross the horizontal axis? Say what each of those two prices means
        for the council --- the reasons are different.

\ansline{At $p=0$ the wristbands are free: plenty go out, no money comes in.}\\
\ansline{At $p=12$ the price is too high: nobody buys, so revenue is \$0 again.}\\

  \item Read the absolute maximum off the graph: which price, and how much revenue? Then check it
        against the council's formula $R(p)=-5p^2+60p$.
\begin{work}
  R(6) &= -5(6)^2+60(6) \\
       &= -180+360 \\
       &= 180
\end{work}

  \item On what interval is revenue increasing? Decreasing? What are the domain and the range in
        this situation?

\ansline{Increasing for $0<p<6$, decreasing for $6<p<12$.}\\
\ansline{Domain $0 \le p \le 12$ dollars; range $0 \le R \le 180$ dollars.}\\
\end{enumerate}
\end{practicebox}

\end{document}
